\subsection{Open-World Planning and Learning}
\label{sec:rw_planning_learning}

Planning and learning provide the reasoning and decision-making capabilities that allow agents to select actions in pursuit of goals. Classical symbolic planning, formalized through STRIPS \citep{fikes1971strips} and PDDL \citep{mcdermott1998pddl}, enables structured long-horizon reasoning by searching over sequences of operators with well-defined preconditions and effects. Modern planners such as Fast Downward \citep{helmert2006fast} achieve strong performance across diverse domains through domain-independent heuristics. However, symbolic planning assumes that the domain model is complete and correct; when the model is inaccurate or missing operators, generated plans may fail during execution.

Reinforcement learning offers a complementary approach, enabling agents to learn effective behaviors through interaction without requiring explicit models. Deep RL algorithms, including DQN \citep{mnih2015human}, PPO \citep{schulman2017proximal}, and SAC \citep{haarnoja2018soft}, have demonstrated success in complex control tasks. Model-based approaches, including world models such as Dreamer \citep{hafner2023mastering} and planning-based methods such as MuZero \citep{schrittwieser2020mastering}, learn environment dynamics and use them for decision-making. These methods reduce the need for hand-specified models but typically assume that the learned dynamics remain stationary.

A recent and rapidly growing line of work leverages large pretrained models for robotic planning and decision-making. Vision-language-action (VLA) models such as RT-2 \citep{pmlr-v229-zitkovich23a} directly map visual observations and language instructions to robot actions, drawing on the broad world knowledge encoded in large language and vision models. SayCan \citep{ahn2022can} and related approaches \citep{huang2022language} use large language models as high-level planners that decompose natural language instructions into sequences of learned skills. Generalist robot policies such as Octo \citep{team2024octo} and $\pi_0$ \citep{black2024pi_0} aim to learn cross-embodiment policies from large-scale multi-robot datasets. While these approaches demonstrate impressive breadth, they face challenges related to physical grounding, hallucination of infeasible actions, and robustness to environmental conditions that diverge from their training data. In open-world settings, where the nature of the change may be qualitatively different from anything in the training distribution, the reliability of these systems remains an open question.
